\section{Images and preimages answer different questions}\label{ch03:sec:images}
Let $f:A\to B$. For $S\subseteq A$, the \emph{image} of $S$ is
\[
f(S)=\{f(x):x\in S\},
\]
the set of outputs of inputs from $S$. It answers ``where do these inputs go?'' When $S=A$, this is the image $f(A)$ from Section~\ref{ch03:sec:functions}. Going the other way, for $T\subseteq B$, the \emph{preimage} of $T$ is
\[
f^{-1}(T)=\{x\in A:f(x)\in T\},
\]
the set of inputs whose outputs land in $T$. It answers ``which inputs end up here?'' Despite the notation, $f^{-1}(T)$ makes sense for every function $f$, whether or not $f$ has an inverse. When $T=\{b\}$, the preimage $f^{-1}(\{b\})$ is the set of all preimages of $b$ in the sense of Section~\ref{ch03:sec:injsurj}. So $f$ is surjective exactly when every such set has at least one member, and injective exactly when every such set has at most one member.

For the square map $f:\RR\to\RR$, $f(x)=x^2$, here are some images and preimages:
\begin{gather*}
f(\{-1,1,2\})=\{1,4\},\qquad f^{-1}(\{1\})=\{-1,1\},\\
f^{-1}(\{4,9\})=\{-3,-2,2,3\},\qquad f^{-1}(\{-1\})=\varnothing.
\end{gather*}
Notice the braces: the value $f(3)=9$ is a number, while the image $f(\{3\})=\{9\}$ is a set. Images and preimages are always sets. The set $T$ may also be an interval: $f^{-1}([0,4])=[-2,2]$, because $0\le x^2$ holds for every real $x$, and $x^2\le4$ holds exactly when $-2\le x\le2$ (Section~\ref{ch03:sec:sets}).

The square map has no inverse, yet all of its preimages make sense. When $f$ does have an inverse $h:B\to A$, the two readings of $f^{-1}(T)$ agree: the preimage of $T$ under $f$ is the image $h(T)$ of $T$ under the inverse. Indeed, take $x\in A$. If $x$ lies in the preimage, then $t=f(x)$ belongs to $T$, and $x=h(f(x))=h(t)$, so $x\in h(T)$. Conversely, if $x=h(t)$ for some $t\in T$, then $f(x)=f(h(t))=t\in T$, so $x$ lies in the preimage.

\pauseq{For the square map $f:\RR\to\RR$, compute $f^{-1}(f(\{2\}))$. Do you get back $\{2\}$?}{No: $f(\{2\})=\{4\}$, and $f^{-1}(\{4\})=\{-2,2\}$. Going forward and then backward returns every input that shares an output with $2$, and $-2$ is such an input. In general $S\subseteq f^{-1}(f(S))$, because each $x\in S$ has its output $f(x)$ in $f(S)$. If $f$ is injective, nothing extra comes back: if $f(x)\in f(S)$, then $f(x)=f(s)$ for some $s\in S$, and injectivity gives $x=s\in S$. So $f^{-1}(f(S))=S$ whenever $f$ is injective.}

Preimages fit well with the set operations of Section~\ref{ch03:sec:sets}. For any $S,T\subseteq B$,
\[
f^{-1}(S\cap T)=f^{-1}(S)\cap f^{-1}(T).
\]
We can justify this one membership condition at a time. For any $x\in A$,
\begin{align*}
 x\in f^{-1}(S\cap T)
 &\Longleftrightarrow f(x)\in S\cap T &&\text{definition of preimage}\\
 &\Longleftrightarrow f(x)\in S\text{ and }f(x)\in T
 &&\text{definition of intersection}\\
 &\Longleftrightarrow x\in f^{-1}(S)\text{ and }x\in f^{-1}(T)
 &&\text{definition of preimage}\\
 &\Longleftrightarrow x\in f^{-1}(S)\cap f^{-1}(T)
 &&\text{definition of intersection}.
\end{align*}
Each step is an ``if and only if,'' valid in both directions. Reading the chain downward shows that every member of the left side lies in the right side, and reading it upward shows the reverse containment. Both sides are subsets of $A$, and $x\in A$ was arbitrary, so the two sets are equal. No inverse function was used. (Exercise~\ref{ex:3.6} asks whether images behave as well.)

\begin{principle}{Habit 3: identify the information actually used}
When you read or write a proof, ask which facts it uses and where. A statement about every target can be much stronger than one about the outputs that are reached. In Proposition~\ref{thm:compconverse}, the composition $g\circ f$ tests $g$ only on the image $f(A)$, so a conclusion about $g$ on all of $B$ asks for more than the argument can give, and here more than is true. The same question can show that a result reaches further than it seems: the proof of Proposition~\ref{ch03:prop:distributive} used only the meanings of ``and'' and ``or'' applied to membership, never what the elements were, so it holds for sets of numbers, people, tasks, or any other objects. Before you propose a generalization, or accept one from an assistant, find the part of each object that the proof actually visits.
\end{principle}
